\chapter{Kapcsolódó munkák}
\label{ch:related_work}

Ebben a fejezetben összefoglaljuk azokat a korábbi kutatásokat és formalizációkat, amelyek közvetlenül kapcsolódnak az elsőrendű logika algebrai szemléletű megközelítéséhez és Agdában történő formalizálásához. Áttekintjük a témához szorosan kapcsolódó dolgozatokat, beszámolókat és egyéb formális rendszerekben végzett hasonló munkákat.

A kapcsolódó munkák között érdemes megemlíteni Fenyvesi Róbert MSc disszertációját \cite{fenyvesi2022}, amely a nulladrendű klasszikus logika különböző megadásainak ekvivalenciáját vizsgálja algebrai úton, az Agda segítségével. Munkája közeli rokonságot mutat jelen dolgozat céljaival, különösen az algebrai szemlélet miatt, azonban a nulladrendű logikával foglalkozik több oldalról, míg a jelen dolgozat az elsőrendű logikát célozza meg.

Samy Avrillon beszámolója \cite{avrillon2023}, amely az ELTE-n töltött gyakornoki kutatási eredményeit foglalja össze, jelen dolgozat elődjének tekinthető, ugyanis a célkitűzések és az alkalmazott módszerek nagyban hasonlók. A munkájában részletesebben tárgyalja az elsőrendű logika formalizálásának módszertani alapjait, különös tekintettel a másodrendű általánosított algebrai elméletek (SOGAT) általánosított algebrai elméletekké (GAT) történő átírására támaszkodva. A diplomamunka célja e megközelítés továbbfejlesztése, a formalizálás teljessé tétele, többféle modell (Bool, Tarski, Family) definiálása, valamint az iniciális modell és a szemantikák részletes kidolgozása, ezek összehasonlítása és további tulajdonságok vizsgálata.

Altenkirch, Hofmann és Streicher klasszikus munkája \cite{altenkirch1995} a nulladrendű intuicionista logika teljességének bizonyítására koncentrál, kategóriaelméleti eszközökkel.

Az algebrai szemlélet elméleti hátteréül szolgál többek között Awodey és Bauer bevezetője a kategóriaelméleti logikába \cite{awodeybauer2020}. Bár ez a megközelítés nem gépileg formalizált, a logika algebrai megközelítéséhez szoros kapcsolat fűzi. Átfogó képet nyújt az algebrai struktúrák és a logikai rendszerek kapcsolatáról.

A témához szorosan kapcsolódik több Coq-ban végzett formalizálás is, amelyek közül az egyik \cite{coq1st} az elsőrendű logika formális leírására és annak metateóriájára koncentrál, elsősorban gyakorlati, típuselméleti eszközöket alkalmazva. Bár nem algebrai megközelítést követ, a formalizálási technikák hasonlóságai miatt hasznos összehasonlítási alapot kínál.

Egy másik kapcsolódó munka \cite{coq0th} a nulladrendű logika teljességi tételének formalizált bizonyítását nyújtja Coq-ban. Ez a munka különösen fontos a jelen dolgozat szempontjából, mivel ugyanarra a logikai fragmentumra koncentrál, mint a dolgozat algebrai irányból történő megközelítése, ám eltérő eszközökkel dolgozik.